\metatitle{Rational shuffle and Delta conjectures}

\metadescription{Rational shuffle, Hikita polynomials, square-path formulas,
Delta conjectures, and related nabla conjectures.}
\metakeywords{rational shuffle theorem, Hikita polynomial, square paths,
Delta conjecture, nabla operator, parking functions}


\section[rationalShuffle]{Rational Shuffle theorem}

A bigraded deformation of the $q,t$-Catalan numbers is introduced in \cite{GorskyMazin2013}.
The rational shuffle conjecture is then stated in \cite[Eq.~(51)]{GorskyNegut2015}
and independently in \cite{Hikita2014} and \cite{Armstrong2012}.

The \defin{extended rational shuffle conjecture}
\label{extendedRationalShuffleConjecture} is first stated in
\cite{BergeronGarsiaLevenXin2015}.
They introduce certain operators $Q_{m,n}$, generalizing operators
$\widetilde{P}_{m,n}$ by Gorsky and Mazin (which required that $m,n$ be coprime),
and conjecture a combinatorial formula for the operators when applied to $(-1)^n$.

Note that $Q_{n+1,n} (-1)^n = \nabla \elementaryE_n$,
so we have the shuffle theorem as a special case.


\section[hikita]{Hikita polynomials}

\begin{polydata}{hikita}
  Name     & Hikita polynomials \\
  Space    & Sym \\
  Basis    & False \\
  Rating   & 1 \\
  Bib      & Hikita2014 \\
  Year     & 2014 \\
  Symbol   & $H_{m,n}(\xvec;q,t)$ \\
  Keywords & fillings, schur-positive, LLT, Frobenius \\
  Category & Schur \\
  PositiveIn & gessel | Hikita2014 | semiring=NN[q,t] \\
  PositiveIn & LLT | BergeronGarsiaLevenXin2015 | scope=vertical-strip LLT polynomials \\
\end{polydata}

\subsection[hikitaPolynomialDefinition]{Definition and extensions}

Using notation similar to that of \cite[Eq.~(11.97)]{Haglund2015},
and \cite{QiuRemmel2018x}, we set, for coprime $m,n$,
\[
H_{(m,n)}(\xvec;q,t) \coloneqq
\sum_{(m,n) \text{ parking functions } P} q^{\area(P)} t^{\dinv(P)} \gessel_{n,\DES(P)}(\xvec).
\]
The functions $H_{(m,n)}(\xvec;q,t)$ are also known as
\defin{Hikita polynomials}\label{hikitaPolynomial} (for coprime $m,n$),
after the introduction in \cite{Hikita2014}.
\name{François Bergeron}, \name{Adriano Garsia}, \name{Emily Leven}, and
\name{Guoce Xin} then introduce the
\defin{extended Hikita polynomials}\label{extendedHikitaPolynomial},
\[
H_{(km,kn)}(\xvec;q,t) \coloneqq
\sum_{(km,kn) \text{ parking functions } P} [ret(P)]_{1/t}
q^{\area(P)} t^{\dinv(P)} \gessel_{n,\DES(P)}(\xvec),
\]
where \defin{$ret(P)$}\label{retStatistic} is the smallest integer $j \gt 0$ such
that the supporting path of $P$ goes through $(jm,jn)$.
They show \cite[Eq.~(6.3)]{BergeronGarsiaLevenXin2015} that
$H_{(m,n)}$ is a positive linear combination of vertical-strip
\hyperref[lltDefinitionTuples]{LLT polynomials}.
\name{Ryan Kaliszewski} and \name{Debdut Karmakar} study a close relationship
between Hikita polynomials and
\hyperref[catalanSymmetric]{Catalan symmetric functions}
\cite{KaliszewskiKarmakar2018}.

\name[Nicolle Gonzalez]{Nicolle González}, \name[Jose Simental]{José Simental},
and \name{Monica Vazirani} introduce higher-rank $(q,t)$-Catalan polynomials
\cite{GonzalezSimentalVazirani2023x}.  These are finite-variable truncations of
Hikita polynomials and are related to parabolic affine
\hyperref[springerFibers]{Springer fibers}.  Their
finite rational shuffle theorem uses rank-$r$ semistandard
\hyperref[parkingFunctions]{parking functions},
with a dinv statistic coming from standardization and a codinv statistic
matching the dimensions of cells in an affine paving.

\name{Alessandro Iraci}, \name{Roberto Pagaria}, \name{Giovanni Paolini}, and
\name{Anna Vanden Wyngaerd} formulate rectangular analogues of the square
paths conjecture and the rise version of the Delta conjecture
\cite{IraciPagariaPaoliniVandenWyngaerd2023}.  They prove their rectangular
paths conjecture when the side lengths of the rectangle are coprime.

\name{Andy Wilson} defines symmetric functions $L_{M,N}$ as weighted sums over
tuples of lattice paths and shows that they satisfy a recursion lifting the
triply graded Khovanov--Rozansky homology of torus links
\cite{Wilson2023TorusLink}.  This gives the torus-link homology as a
specialization of $L_{M,N}$ and relates the construction conjecturally to
elliptic Hall algebra operators.


\subsection[hikitaOperators]{The operators $Q_{m,n}$}

The operators $Q_{m,n}$ are defined recursively via the
Bergeron--Garsia operators \defin{$D_k$}\label{bergeronGarsiaOperator}
(acting on $F \in \spaceSym$)
\[
D_k F[X] \coloneqq [z^k] F\left[ X + \frac{(1-q)(1-t)}{z} \right] \sum_{j \geq 0}
(-z)^j \elementaryE_j[X].
\]

The construction of the $Q_{m,n}$ is involved, and different normalizations
in the literature can introduce sign differences.

\name{Eugene Gorsky} and \name{Andrei Neguț} prove the first identity
below for coprime $m,n$ \cite{GorskyNegut2015}; \name{Dun Qiu} and
\name{Jeffrey Remmel} prove the second \cite{QiuRemmel2018x}:
\[
 \nabla Q_{m,n} \nabla^{-1} = Q_{m+n,n} \qquad
 \nabla Q_{kn,n} \nabla^{-1} = Q_{(k+1)n,n}.
\]

\name{Viviane Pons} gives an explicit description of the zeta map on
\hyperref[qtCatalan]{Dyck paths} in terms of area sequences
\cite{Pons2022x}.  This provides a direct and implementable way to see how
the zeta map sends the pair of statistics
$(\dinv,\area)$ to $(\area,\bounce)$.


\subsection[rationalShuffleTheorem]{The rational shuffle theorem}

\name{Anton Mellit} proved the following theorem \cite{Mellit2016x}.
\begin{theorem}[The (extended) rational shuffle theorem]
For coprime $m,n$ and $k \geq 0$, we have
\[
 H_{(km,kn)}(\xvec;q,t) =  Q_{km,kn} (-1)^n.
\]
\end{theorem}
In fact, he proves the compositional refinement
stated in \cite[Conj.~3.3]{BergeronGarsiaLevenXin2015}.


\subsection[hikitaGeometrySpecializations]{Geometry and specializations}

\name{Tatsuyuki Hikita} shows that $B_{m,n}(\xvec;q,t)$
is the Frobenius character of the action in the homology of a
\hyperref[springerFibers]{Springer fiber} in the affine
\hyperref[flagVarieties]{flag variety} equipped
with a specified filtration \cite{Hikita2014}.


The \hyperref[prelimQanalogsCatalan]{rational Catalan numbers}
are defined as $\catalan_{a/b} = \frac{1}{a+b}\binom{a+b}{a,b}$ for coprime $a,b$.
For coprime $m,n$,
\[
 \langle H_{(m,n)}(\xvec;q,t), \elementaryE_n \rangle_* = \catalan_{m/n}.
\]
Here $\langle \cdot, \cdot \rangle_*$ is a deformation of the
\hyperref[hallInnerProduct]{Hall inner product}, see
\cite[Eq.~(4.6) and Eq.~(6.29)]{BergeronGarsiaLevenXin2015}.

There is also an expression for $H_{(m,n)}(\xvec;q,t)$ using
\defin{Tesler matrices}\label{teslerMatrix}.


See \cite{QiuRemmel2018x} for results
on Schur coefficients of $H_{(m,n)}(\xvec;q,t)$.

\name{Maria Gillespie}, \name{Eugene Gorsky}, and \name{Sean T. Griffin}
show that the symmetric function in the Rise Delta theorem can be obtained
from the Rational Shuffle theorem by applying a
\hyperref[schurSkew]{Schur skewing operator}
\cite{GillespieGorskyGriffin2024x}.  Their combinatorial proof gives a new
derivation of the Rise Delta theorem from the Rational Shuffle theorem.



\section[tamariInterval]{Tamari intervals and more sets of variables}

\name{Nantel Bergeron}, \name{Cesar Ceballos}, and \name{Vincent Pilaud}
connect Hopf algebras of \hyperref[pipeDream]{pipe dreams} with multivariate diagonal harmonics
\cite[Sec.~8]{BergeronCeballosPilaud2018x}.  Their dominant pipe-dream Hopf
algebra is related to $\nu$-Tamari lattices, and this leads to the notion of a
\defin{Hopf chain}\label{hopfChain}: a nested tuple of
\hyperref[qtCatalan]{Dyck paths}
$(\pi_1,\dotsc,\pi_r)$ with $\pi_1$ equal to the diagonal path and with every
triple satisfying a compatible $\nu$-Tamari interval condition.

For $n\leq 4$ and any number $r$ of variable sets, they show that the
$q,t$-\hyperref[frobeniusCharacteristic]{Frobenius characteristic} of the multivariate diagonal harmonic space
$\mathcal{DH}_{n,r}$ is given by a sum over Hopf chains, weighted by a collar
statistic and by \hyperref[LLT]{LLT polynomials}.  When $r=3$, the construction specializes to
the Tamari-interval conjectural picture of \name{François Bergeron} and
\name[Louis-Francois Preville-Ratelle]{Louis-François Préville-Ratelle}, where
the length of a longest chain in a Tamari interval is one of the expected
statistics.


\section[dhSquare]{The square paths theorem, $\nabla \powerSum_n$}

The square path conjecture was formulated by \name{Nicholas Loehr} and
\name{Gregory Warrington} \cite{LoehrWarrington2007}.
It provides a combinatorial expression for $\nabla \powerSum_n$
as a sum over preference functions.

\subsection[squarePathPreferenceFunctions]{Preference functions and statistics}

A \defin{preference function}\label{preferenceFunction} is a map
$f:[n] \to [n]$.
A \hyperref[parkingFunctions]{parking function} is any preference function such
that $|f^{-1}([k])| \geq k$ for all $k \in [n]$.
With the parking language, we think of $f(i)$ as the
\defin{preference}\label{parkingPreference} of car $i$.
Associate a lattice path to a preference function as follows.
In an $n\times n$-square, write the numbers $f^{-1}(i)$ in column $i$,
starting from \(i=1\).
The numbers in each column are increasing upwards, and the $j$th number written
is placed in row $j$.
Then there is a unique lattice path from $(0,0)$ to $(n,n)$ using North and East
steps, such that the numbers are written immediately to the right of the North
steps.
See \hyperref[preferenceFunctions]{preference functions and square paths}
for a more detailed definition of the reading word and dinv statistic.

\begin{example*}[Path from preference function]
For example, the
preference function $(1,5,1,2,1)$ gives rise to the path
\(P=\mathtt{nnneneeene}\),
and the diagram
\begin{figure}
\begin{ytableau}
 & & & & 2 \\
 & 4 & & &   \\
5 & & & &  \\
3 & & & &  \\
1 & & & &
\end{ytableau}
\end{figure}
\end{example*}


The \defin{area}\label{preferenceFunctionArea} of a preference function is the
number of full cells above the
lowest diagonal which contains a car. Thus, if the path is a
\hyperref[qtCatalan]{Dyck path} (and the preference function is a parking function),
the area coincides with the classical notion of area.
The \defin{dinv}\label{preferenceFunctionDinv} statistic is similar to the
inv-statistic for \hyperref[LLT]{LLT polynomials},
but one must add the number of cars strictly below the $x=y$ diagonal.

\subsection[squarePathTheorem]{The square paths theorem}

\name{Emily Sergel} gave a proof of the square path conjecture \cite{Sergel2017}.
\begin{theorem}[Sergel (2017)]
We have the expansion in the \hyperref[gessel]{fundamental quasisymmetric functions}
\[
(-1)^{n-1} \nabla \powerSum_n = \sum_{w \in Pref(n)} t^{\area(w)} q^{\dinv(w)} \gessel_{\mathrm{IDES}(w)}(\xvec),
\]
where the sum is taken over all preference functions of size $n$.
\end{theorem}

\subsection[squarePathSpecializations]{Narayana coefficients and specializations}

Since \(\omega(\powerSum_n)=(-1)^{n-1}\powerSum_n\), the left-hand side is
\(\nabla\omega(\powerSum_n)\).  It is tempting to call some of the
coefficients below type \(B\) Narayana numbers, because the specialization
at \(q=t=1\) gives the classical type \(B\) Narayana numbers
\(\binom{n}{k}^2\).  This should be understood only as an analogy:
the finite-reflection-group type \(B\) Catalan story has different
combinatorics, involving shifted or reflection-group paths
\cite{Stump2010FussCatalan,Lu2019TypeBCatalan}.

\begin{proposition}[Square-path Narayana coefficients at \(q=t=1\)]
For \(0\leq k\leq n\),
\[
\left.
\langle \nabla\omega(\powerSum_n),
\completeH_k\completeH_{n-k}\rangle
\right|_{q=t=1}
=
\binom{n}{k}^2.
\]
\end{proposition}
\begin{proof*}
By Sergel's theorem, the word version of the expansion, and Hall duality,
the scalar product at \(q=t=1\) counts
\hyperref[wordPreferenceFunctionDefinition]{word preference functions} of size
\(n\) with monomial \(x_1^kx_2^{n-k}\).
With only labels \(1\) and \(2\), each of the \(n\) columns is one of
\[
\emptyset,\qquad \{1\},\qquad \{2\},\qquad \{1,2\}.
\]
Suppose there are \(j\) columns of type \(\{1,2\}\).  Then there are
\(k-j\) columns of type \(\{1\}\), \(n-k-j\) columns of type \(\{2\}\),
and \(j\) empty columns.  Hence the number is
\[
\sum_{j\geq 0}
\binom{n}{j}\binom{n-j}{n-2j}\binom{n-2j}{k-j}.
\]
This simplifies to
\[
\binom{n}{k}\sum_{j\geq0}\binom{k}{j}\binom{n-k}{j}
=\binom{n}{k}^2
\]
by Vandermonde's identity.
\end{proof*}

\begin{proposition}[Specialization at \(t=1/q\)]
Let
\[
  S_{n,k}(q,t)\coloneqq
  \langle \nabla\omega(\powerSum_n),
  \completeH_k\completeH_{n-k}\rangle.
\]
Then
\[
  q^{k(n-k)}S_{n,k}(q,q^{-1})
  =
  \qbinom{n}{k}_q^2.
\]
\end{proposition}
\begin{proof*}
This is a coefficient extraction from
\cite[Thm.~4.30]{DAdderioIraciWyngaerd2022}.  Since
\(\nabla=\Delta_{\elementaryE_n}\) on homogeneous symmetric functions of
degree \(n\), their theorem gives
\[
\left.\nabla\omega(\powerSum_n)\right|_{t=1/q}
=
q^{-\binom{n}{2}}\elementaryE_n[X[n]_q].
\]
The coefficient dual to \(\completeH_k\completeH_{n-k}\) is the coefficient of
the monomial \(x_1^kx_2^{n-k}\).  Choosing the \(k\) \(q\)-weights attached to
\(x_1\) and the \(n-k\) \(q\)-weights attached to \(x_2\) gives
\[
\left[x_1^kx_2^{n-k}\right]
q^{-\binom{n}{2}}\elementaryE_n[X[n]_q]
=
q^{-k(n-k)}\qbinom{n}{k}_q^2.
\]
\end{proof*}

The same result of D'Adderio--Iraci--Vanden Wyngaerd gives the sign-character
specialization
\[
\left.
\langle \nabla\omega(\powerSum_n),\elementaryE_n\rangle
\right|_{t=1/q}
=
q^{-\binom{n}{2}}\frac{1}{1+q^n}\qbinom{2n}{n}_q,
\]
see \cite[Prop.~8.3]{DAdderioIraciWyngaerd2022}.

\subsection[squarePathLLTSchedules]{LLT expansions and schedules}

This expression can also be written as a sum over vertical-strip
\hyperref[lltDefinitionTuples]{LLT polynomials}.
A path arising from a preference function of size $n$ can be described using a weak
composition $\alpha$, such that $\alpha_i \coloneqq |f^{-1}(i)|$.
The path determines $\area(\alpha)$
and $\mathrm{carsBelow}(\alpha)$, the number of cars below the main diagonal.
Finally, we can associate an $n$-tuple of skew shapes, $\nuvec$,
determined by the positions of the cars. This gives the following:
\begin{proposition}
\begin{equation*}
(-1)^{n-1} \nabla \powerSum_n
= \sum_{\alpha} t^{\area(\alpha)} q^{\mathrm{carsBelow}(\alpha)}
\LLT_{\nuvec(\alpha)}(\xvec;q)
\end{equation*}
where the sum ranges over all \hyperref[weakCompositions]{weak compositions} of length $n$ and size $n$.
\end{proposition}


\begin{example*}[Schur-expansion of $(-1)^{n-1} \nabla \powerSum_n$]
\begin{array}{ll}
\toprule
n & (-1)^{n-1} \nabla \powerSum_n \\
\midrule
1 & \schurS_{1} \\
2 & \schurS_{2} + (q + t + q t) \schurS_{11} \\
3 & \schurS_{3} +  (q + q^2 + t + q t + q^2 t + t^2 + q t^2 + q^2 t^2)\schurS_{21} + \\
  & +(q^3 + q t + q^2 t + q^3 t + q t^2 + q^2 t^2 + q^3 t^2 + t^3 + q t^3 +
   q^2 t^3) \schurS_{111} \\
\bottomrule
\end{array}
\end{example*}

\name{Emily Sergel}'s proof uses \defin{schedules}\label{schedule} to enumerate preference
functions with a fixed diagonal word.  If $\tau$ has run lengths
$\rho_k,\dotsc,\rho_0$ and $w^{(l)}(c)$ denotes the $l$-schedule number of a
car $c$, then \cite[Thm.~2.1]{Sergel2017} gives the product formula
\[
\sum_{\substack{\operatorname{diagword}(Pr)=\tau\\ l(Pr)=l}}
  t^{\area(Pr)}q^{\dinv(Pr)}
= t^{\operatorname{maj}(\tau)}
  q^{\rho_0+\dotsb+\rho_{l-1}}\prod_{c=1}^n [w^{(l)}(c)]_q.
\]
The key shift step is that the multiset of $l$-schedule numbers differs from
the ordinary parking-function schedule by replacing one run length:
\[
  \{w^{(l)}(c):1\leq c\leq n\}
  =
  \{w_i:1\leq i\leq n\}\cup\{\rho_l\}\setminus\{\rho_0\}.
\]
This is \cite[Thm.~3.1]{Sergel2017}, and its proof uses
\cite[Lem.~3.1]{Sergel2017}, a conjugate-partition multiset identity.




\section[dhDelta]{The Delta conjecture, $\Delta'_{\elementaryE_k} \elementaryE_n$}

The Delta conjecture generalizes the Shuffle theorem, and was stated by
\name{James Haglund}, \name{Jeffrey Remmel}, \name{Andrew Wilson}
\cite{HaglundRemmelWilson2018}.  It has two versions, the
\defin{rise version}\label{riseDeltaConjecture} and the
\defin{valley version}\label{valleyDeltaConjecture}.  The rise version is now
a theorem: it follows
from the compositional Delta theorem of \name{Michele D'Adderio} and
\name{Anton Mellit} \cite{DAdderioMellit2022}, and also from the extended Delta
theorem of \name{Jonah Blasiak}, \name{Mark Haiman}, \name{Jennifer Morse},
\name{Anna Pun}, and \name{George H. Seelinger}
\cite{BlasiakHaimanMorsePunSeelinger2023}.

\begin{theorem}[Rise Delta theorem]
\[
\Delta'_{\elementaryE_k} \elementaryE_n
= [z^{n-k-1}]
\sum_{w \in WPF(n)}
  t^{\area(w)}q^{\dinv(w)} \xvec_w
\prod_{i \in Rise(w)} \left(1 + z/t^{a_i(w)}\right)
\]
\end{theorem}

The valley version remains open in full generality; in particular, symmetry of
the combinatorial side is still part of the problem.
\name{James Haglund} and \name{Emily Sergel} introduce schedule formulas for
the valley side and use them to prove several special cases
\cite{HaglundSergel2020}.
\begin{conjecture}[Valley Delta conjecture, J. Haglund, J. Remmel, A. Wilson]
\[
\Delta'_{\elementaryE_k} \elementaryE_n =
[z^{n-k-1}]
\sum_{w \in WPF(n)} t^{\area(w)}q^{\dinv(w)} \xvec_w
\prod_{i \in Val(w)} \left(1 + z/t^{d_i(w)+1}\right)
\]
\end{conjecture}

The rise version has a compositional generalization, which is the result proved
in \cite{DAdderioMellit2022}.

\name{Dun Qiu} and \name{Minhao Zhang} prove a nonsymmetric compositional
Delta theorem using flagged \hyperref[LLT]{LLT polynomials}
\cite{QiuZhang2026xNonsymmetricDelta}.  Their signed and unsigned identities
use nonsymmetric analogues of $\nabla$ and $\tau^*$; Weyl symmetrization
recovers the ordinary compositional Delta theorem.  They also formulate
stable-atom positivity conjectures for the nonsymmetric identities.

\name{Alessandro Iraci} and \name{Anna Vanden Wyngaerd} use a pushing
operation to pass from the valley Delta conjecture to the generalized valley
Delta conjecture \cite{IraciWyngaerd2022}.

\name{Michele D'Adderio}, \name{Alessandro Iraci}, \name{Yvan Le Borgne},
\name{Marino Romero}, and \name{Anna Vanden Wyngaerd} use Theta operators to
give a symmetric-function formula enumerating tiered trees
\cite{DAdderioIraciLeBorgneRomeroVandenWyngaerd2023}.  They also formulate a
general conjecture whose special cases are connected to unified forms of the
Delta conjecture.

\name{Abel Lacabanne}, \name{Pedro Vaz}, and \name{Arik Wilbert} study
two-row \hyperref[springerFibers]{Delta Springer varieties}, introduced by Griffin--Levinson--Woo to
geometrically realize a symmetric function from the Delta conjecture
\cite{LacabanneVazWilbert2025}.  They describe the irreducible components
combinatorially, compare them with ordinary and exotic Springer fibers, and
extend the homology action to a degenerate affine Hecke algebra action.
\name{Sean T. Griffin}, \name{Jake Levinson}, and \name{Alexander Woo}
introduced these Delta Springer fibers and proved a $t=0$ Delta-conjecture
formula through their geometry \cite{GriffinLevinsonWoo2024}.

\name{Alessandro Iraci}, \name{Philippe Nadeau}, and
\name{Anna Vanden Wyngaerd} prove a unified Delta theorem at $t=0$ using
segmented Smirnov words \cite{IraciNadeauWyngaerd2023x}.  The symmetric
function they interpret is related to a diagonal coinvariant ring with one set
of commuting variables and two sets of anticommuting variables.

\name{Sylvie Corteel}, \name[Matthieu Josuat-Verges]{Matthieu Josuat-Vergès},
and \name{Anna Vanden Wyngaerd} prove a specialization of the Delta conjecture
at $q=-1$ \cite{CorteelJosuatVergesWyngaerd2024}.  The result gives a
$t$-analogue of Euler numbers in terms of peaks and valleys of permutations.

\name{Michele D'Adderio} and \name{Alessandro Iraci} prove several consequences
of the valley Delta conjectures \cite{DAdderioIraci2023}, including the
Schröder case of the valley Delta conjecture and the Schröder case of its
square version.


\subsection[dhGeneralizedDelta]{The generalized Delta conjecture,
$\Delta_{\completeH_m} \Delta'_{\elementaryE_{n-k-1}} \elementaryE_n$
}

The generalized valley Delta conjecture is equivalent to the original valley
Delta conjecture, by work of \name{Alessandro Iraci} and
\name{Anna Vanden Wyngaerd} \cite{IraciWyngaerd2022}.  Their work also relates
the valley Delta conjecture to the corresponding valley Delta square
conjectures.

\subsection[dhExtendedDelta]{The extended Delta conjecture, $\Delta'_{\elementaryE_k} \Delta_{\completeH_r} \elementaryE_n$}

\name{James Haglund}, \name{Jeffrey Remmel}, \name{Andrew Wilson}
\cite{HaglundRemmelWilson2018}
also formulated the \defin{extended Delta conjecture}
\label{extendedDeltaConjecture},
which gives a combinatorial formula for
\[
\Delta'_{\elementaryE_k} \Delta_{\completeH_r} \elementaryE_n, \quad k \lt n.
\]

\begin{theorem}[Extended Delta theorem]
\[
\Delta'_{\elementaryE_k} \Delta_{\completeH_r} \elementaryE_n
= [z^{n-k-1}]
\sum_{w \in WPF(n,r)}
  t^{\area(w)}q^{\dinv(w)} \xvec_w
\prod_{i \in Rise(w)} \left(1 + z/t^{a_i(w)}\right)
\]
\end{theorem}
This was proved in \cite{BlasiakHaimanMorsePunSeelinger2023}.

\name{Dun Qiu} and \name{Andrew Wilson} \cite{QiuWilson2020} formulated the
valley version of the extended Delta conjecture.
\begin{conjecture}[Valley extended Delta conjecture]
\[
\Delta'_{\elementaryE_k} \Delta_{\completeH_r} \elementaryE_n
=
[z^{n-k-1}]
\sum_{w \in WPF(n,r)} t^{\area(w)}q^{\dinv(w)} \xvec_w
\prod_{i \in Val(w)} \left(1 + z/t^{d_i(w)+1}\right)
\]
\end{conjecture}
The valley version has been proved true for $q=0$ or $t=0$.

See also \cite{Bergeron2024}.
For consequences and related special cases, see \cite{DAdderioIraci2023}.


\section[dhDeltaSquare]{Generalized Delta square conjecture}

The \defin{generalized Delta square conjecture}
\label{generalizedDeltaSquareConjecture} (stated in
\cite{DAdderioIraciWyngaerd2019})
is about a combinatorial interpretation of
\[
  \frac{[n-k]_t}{[n]_t} \Delta_{\completeH_m} \Delta_{\elementaryE_{n-k}} \omega(\powerSum_n).
\]
Remark~3.14 of \cite{DAdderioIraciWyngaerd2019} shows that this, if the
conjecture is true, can be expressed as a sum of \hyperref[LLT]{LLT polynomials}.

The full Delta square conjectures remain open in general.  The valley version
was formulated in \cite{IraciWyngaerd2021}, and the implications from the
original valley Delta conjecture are discussed in \cite{IraciWyngaerd2022}.

Some related questions are solved in \cite{Romero2022}.

\name{Alessandro Iraci}, \name{Roberto Pagaria}, and \name{Giovanni Paolini}
prove a fall-decorated rational shuffle theorem
\cite{IraciPagariaPaolini2025x}.  This gives a rational analogue of the
fall Delta theorem and is designed to interact with the Delta square
conjectures.

\name{Sylvie Corteel}, \name{Alexander Lazar}, and
\name{Anna Vanden Wyngaerd} study decorated square-path enumerators at
$q=-1$ \cite{CorteelLazarWyngaerd2024x}.  Using a cyclic cutting-and-pasting
action, they show that the valley Delta square enumerator vanishes when
$n-k$ is even, while in the odd case it becomes a positive polynomial related
to $t$-Euler numbers.

\name{John Lentfer} proposes a monomial basis for the
$(1,2)$-bosonic--fermionic coinvariant ring \cite{Lentfer2024x}.  The
candidate basis interpolates between the Kim--Rhoades modified Motzkin-path
basis and the super-Artin basis, has Zabrocki's conjectured cardinality
$2^{n-1}n!$, and gives conjectural Hilbert and Frobenius series equivalent
to the segmented-Smirnov-word formulas of Iraci--Nadeau--Vanden Wyngaerd.
\name{Mike Zabrocki} conjectures a module for the Delta conjecture
\cite{Zabrocki2019x}; \name{Nolan Wallach} studies consequences of this
conjectural module for the action of $\symS_n$ on polynomial differential forms
\cite{Wallach2019x}.  See also \name{François Bergeron}'s
bosonic--fermionic diagonal coinvariant modules conjecture
\cite{Bergeron2020BosonicFermionic}.


\section[dhSchur]{Conjecture on $\nabla^m \schurS_\lambda$}

\name{Nicholas Loehr} and \name{Gregory Warrington} conjectured
a combinatorial expansion of $\nabla^m \schurS_\lambda$.

This conjecture (in fact, a much stronger result)
was proved by \name{Jonah Blasiak}, \name{Mark Haiman}, \name{Jennifer Morse},
\name{Anna Pun}, and \name{George H. Seelinger},
in \cite{BlasiakHaimanMorsePunSeelinger2021xLWConj,BlasiakHaimanMorsePunSeelinger2021xLLT}.
Their work also gives a new proof of the Shuffle theorem and similar
identities.
\name{Donghyun Kim} and \name{Jaeseong Oh} introduce Macdonald piece
polynomials and obtain the Loehr--Warrington conjecture as a corollary of a
more general framework \cite{KimOh2024x}.
\name{François Bergeron} studies related $\GL_k\times\symS_n$-modules and
nabla applied to hook-indexed Schur functions \cite{Bergeron2019HooksNabla}.


\section[dhMonomial]{The monomial nabla conjecture}


\begin{conjecture}[Bergeron--Garsia--Haiman--Tesler (1999), \cite[Conj. IV]{BergeronGarsiaHaimanTesler1999}]
We have for every pair of partitions $\lambda$, $\mu$,
\[
(-1)^{|\lambda|-\length(\lambda)} \langle \nabla \monomial_\lambda, \schurS_\mu \rangle \in \setN[q,t].
\]
\end{conjecture}
This also appears in \cite[Conj.~11]{LoehrWarrington2007}.

\name{Dun Qiu} and \name{Minhao Zhang} proved this conjecture in full
generality, and in fact proved the following stronger statement.
\begin{theorem}[Qiu--Zhang \cite[Thm.~1.2]{QiuZhang2026x}]
For all partitions \(\alpha,\lambda\vdash n\) and all integers \(r\geq 1\),
\[
(-1)^{|\alpha|-\length(\alpha)}
\left\langle \nabla^r \monomial_\alpha,\schurS_\lambda \right\rangle
\in \setN[q,t].
\]
\end{theorem}
Their proof expands \((-1)^{|\alpha|-\length(\alpha)}\monomial_\alpha\)
positively in the functions \(C_\beta(1)\), then applies the compositional
shuffle theorems and \hyperref[lltSchurPositivity]{Schur positivity of LLT
polynomials}.  They also prove an
\(\elementaryE\)-positive analogue after the substitution \(q\mapsto q+1\)
\cite[Thm.~1.3]{QiuZhang2026x}.

\name{Menghao Qu} and \name{Guoce Xin} prove
\hyperref[schurPositiveDefinition]{Schur positivity} of
$(-1)^k\nabla \monomial_{2^k1^\ell}$ and give a parking-function
interpretation \cite{QuXin2025}.  Their proof establishes $C$-positivity for
$(-1)^k\monomial_{2^k1^\ell}$ and then applies the compositional shuffle
theorem.




\section[dhOther]{Other conjectures}


The first Bergeron--Garsia--Haiman--Tesler conjecture in this group concerns
signed \hyperref[schurPositiveDefinition]{Schur positivity} of
$\nabla^m\schurS_\lambda$.  This is now a theorem, following from the proof of
the Loehr--Warrington conjecture discussed above.

\subsection[nablaHallLittlewoodBGHT]{Nabla on modified Hall--Littlewood polynomials}

The next statements were Conjectures~II--III of
Bergeron--Garsia--Haiman--Tesler
\cite[Conjs.~II--III]{BergeronGarsiaHaimanTesler1999}.  The compositional
shuffle theorem proves Conjecture~II.  \name{Feng Gu}, \name{Mingyang Kang},
\name{Dun Qiu}, and \name{Minhao Zhang} prove both parts of Conjecture~III
and extend all three statements to every positive power of $\nabla$.

\begin{theorem}[Gu--Kang--Qiu--Zhang,
\cite[Cor.~18, Thm.~8, and Cor.~9]{GuKangQiuZhang2026x}]
For partitions $\mu,\nu,\lambda\vdash n$ and every integer $r\geq1$, we have
\[
(-1)^{n-\length(\mu)}
\left\langle \nabla^r\macdonaldH_\mu(\xvec;0,t),
\schurS_\lambda\right\rangle
\in\setN[q,t]
\]
and
\[
\left\langle \nabla^r\omega
\macdonaldH_\mu(\xvec;0,t^{-1}),\schurS_\lambda\right\rangle
\in\setN[q,t].
\]
If $\mu\trianglerighteq\nu$ in
\hyperref[dominanceOrder]{dominance order}, then
\[
\left\langle
\nabla^r\omega\macdonaldH_\mu(\xvec;0,t^{-1})
-\nabla^r\omega\macdonaldH_\nu(\xvec;0,t^{-1}),
\schurS_\lambda
\right\rangle
\in\setN[q,t].
\]
\end{theorem}

The dominance proof reduces to covers and uses a single-LLT identity; see the
\hyperref[lltNablaHallLittlewoodCovers]{LLT-polynomial formulation}.
\name{Menghao Qu} previously proved the two-column cases and their extension
to positive powers of $\nabla$ \cite{Qu2026xHallLittlewood}.



In \cite{BergeronGarsiaHaimanTesler1999} and later \cite[Conj.~3.20]{Haiman2002},
the following conjecture is stated:
\begin{conjecture}[Bergeron--Garsia--Haiman--Tesler (1999)]
The expression
\[
\nabla_{\schurS_\nu} \elementaryE_n
\]
is a $(q,t)$-Schur-positive polynomial for all $\nu$ and $n$.
\end{conjecture}
